\documentclass[12pt]{article}
\usepackage[T1]{fontenc}
\usepackage[utf8]{inputenc}
\usepackage{lmodern}
\usepackage{textcomp}
\usepackage{enumitem}
\usepackage{geometry}
\geometry{margin=1in}
\begin{document}
\begin{center}
Answer Key - Political Science - Graduate Level Assessment 7 \
\small (Answers are ordered exactly as in the question paper.)
\end{center}

\section*{Multiple Choice}
\begin{enumerate}[label=\arabic*.]
\item \textbf{Answer:} A
\\ \textit{Options:} A) A strand of neo-isolationist sentiment; B) A strand of internationalist sentiment; C) An expression of American cultural superiority; D) Increased incorporation of public opinion in foreign policy making
\\ \textit{Note:} Correct option: A - A strand of neo-isolationist sentiment
\item \textbf{Answer:} D
\\ \textit{Options:} A) A growing effort by governments to prevent the spread of weapons, especially of WMDs.; B) A growth in volume of the arms trade.; C) A change in the nature of the defence trade, linked to the changing nature of conflict.; D) All of these options.
\\ \textit{Note:} Correct option: D - All of these options.
\item \textbf{Answer:} C
\\ \textit{Options:} A) Barry Buzan and Ole Waever; B) Kenneth Waltz and Hans Morgenthau; C) Karl Marx and Friedrich Engels; D) Adam Smith and Karl Marx
\\ \textit{Note:} Correct option: C - Karl Marx and Friedrich Engels
\item \textbf{Answer:} C
\\ \textit{Options:} A) the United Nations.; B) the National Security Council.; C) the State Department.; D) the National Security Agency.
\\ \textit{Note:} Correct option: C - the State Department.
\item \textbf{Answer:} D
\\ \textit{Options:} A) Leaders have incentives to lie; B) If leaders revealed their programs, they would be more likely to be attacked; C) Leaders will not always grant foreign monitors access to their nuclear programs; D) ALL of the above
\\ \textit{Note:} Correct option: D - ALL of the above
\end{enumerate}

\section*{True / False}
\begin{enumerate}[label=\arabic*.]
\setcounter{enumi}{5}
\item \textbf{Answer:} True
\\ \textit{Options:} A) True; B) False
\\ \textit{Note:} LLM-generated true statement based on MCQ.
\item \textbf{Answer:} False
\\ \textit{Options:} A) True; B) False
\\ \textit{Note:} LLM-generated false statement. Correct answer: 'It encourages liberalization of global trade'.
\item \textbf{Answer:} True
\\ \textit{Options:} A) True; B) False
\\ \textit{Note:} LLM-generated true statement based on MCQ.
\item \textbf{Answer:} True
\\ \textit{Options:} A) True; B) False
\\ \textit{Note:} LLM-generated true statement based on MCQ.
\item \textbf{Answer:} False
\\ \textit{Options:} A) True; B) False
\\ \textit{Note:} LLM-generated false statement. Correct answer: 'Regime type'.
\end{enumerate}

\section*{Long Form Response}
\begin{enumerate}[label=\arabic*.]
\setcounter{enumi}{10}
\item \textbf{Answer:} Broader political considerations about the potential implications of long-term mobilization resulted in a shift from tactical operations to long-term strategy.. Explain why this is the correct answer and provide supporting evidence. Reference key facts or concepts that support this choice.
\\ \textit{Note:} Long-form derived from MCQ; award credit for stating the correct idea and giving supporting reasoning.
\item \textbf{Answer:} The subaltern refers to populations that are marginalised or outside of the hegemonic power structure.. Explain why this is the correct answer and provide supporting evidence. Reference key facts or concepts that support this choice.
\\ \textit{Note:} Long-form derived from MCQ; award credit for stating the correct idea and giving supporting reasoning.
\end{enumerate}

\vfill
\noindent\textit{End of Answer Key}
\end{document}